% !TeX root = Lec08_NK_Baisc.tex

\ifdefined\AdvMacroMain
  \chapter{Lecture 8: Basic New Keynesian}
  \let\ThisNoteEnd\relax
\else
  \documentclass[12pt]{ctexart}
  \newcommand{\ProjectRoot}{..}
  \usepackage[a4paper,margin=0.85in]{geometry}
\usepackage{newpxtext}
\usepackage{amsmath,amssymb,amsthm,mathtools}
\usepackage{newpxmath}
\usepackage{xcolor}
\usepackage{enumitem}
\usepackage{graphicx}
\usepackage{array,tabularx}
\usepackage{float}
\usepackage{hyperref}

\definecolor{noteRed}{RGB}{190,45,90}

\newcommand{\red}[1]{\textcolor{noteRed}{#1}}
\newcommand{\blue}[1]{\textcolor{blue}{#1}}
\newcommand{\purple}[1]{\textcolor{purple}{#1}}

\newcommand{\E}{\mathbb{E}}
\newcommand{\dd}{\mathrm{d}}
\newcommand{\MPN}{\mathrm{MPN}}
\newcommand{\MRS}{\mathrm{MRS}}
\newcommand{\MRT}{\mathrm{MRT}}
\newcommand{\Var}{\operatorname{Var}}
\newcommand{\boxit}[1]{\fbox{\ensuremath{#1}}}
\newcommand{\circled}[1]{\textcircled{\scriptsize #1}}

\newenvironment{tightalign}
  {\[\begin{aligned}}
  {\end{aligned}\]}

\graphicspath{
  {\ProjectRoot/_assets/figures/}
  {\ProjectRoot/_assets/figures/lecture_04/}
}

  \title{Lecture 8: Basic New Keynesian}
  \author{Chengyuan He}
  \date{May 2026}
  \begin{document}
  \maketitle
  \def\ThisNoteEnd{\end{document}}
\fi

\section*{RBC to NK}
\textbf{RBC:} Fluctuations explained by technology shocks  (kydland--Prescott 1982; Prescott 1986)
\[
\text{Use DSGE models as tool: FOCs replacing behavioral equations, rational expectations.}
\]
\begin{enumerate}
  \item Efficient allocation.
  \item Implication: stabilization policies non-necessary.
  \item Equilibrium outcome responds to exogenous shocks.
  \item Limited role of monetary factors, neutrality of monetary policy: Friedman's rule
\end{enumerate}
\red{* Contradiction to:}
\begin{enumerate}[label=\arabic*.]
  \item \red{Common belief: influence on output and employment.}
  \item \red{Empirical challenge: non-neutrality.}
\end{enumerate}
\red{$\Rightarrow$ Irrelevancy to policy evaluation: Normative implication different from policy practice.}

\red{Monetary policy can affect real variables or not: e.g. output, employment, consumption.}
\vspace{-2pt}
\par\smallskip

\par\smallskip
\begin{tabularx}{\linewidth}{@{}>{\raggedright\arraybackslash}X>{\raggedright\arraybackslash}X@{}}
\textbf{RBC} & \textbf{New Keynesian} \\
Agents / household / firms & $\sim$ \\
DSGE & $\sim$ \\
Markets are frictionless and perfectly competitive & Goods market imperfection (monopolistic competition) \\
Price adjust instantaneously & Nominal rigidity (price, wage etc.) \\
Technology shocks \newline \red{Not really: RBC also can introduce other shocks.} & Other shocks and transmission mechanisms
\end{tabularx}\section*{New in NK}

\begin{enumerate}[leftmargin=1.7em,itemsep=0.35em,topsep=0.2em]
    \item Monopolistic competition: $P^m>P^c, Y^m<Y^c$, \red{垄断竞争}
    \hfill \red{inefficient response to shocks.}

    \item Nominal rigidities: staggered price / wage, 名义刚性，价格短期刚性.

    \item Short-run non-neutrality of monetary policy, 短期货币非中性
    $\Rightarrow$ \red{policy analysis \& intervention.}
\end{enumerate}

\par\smallskip

These elements are new to RBC but old to Keynesian (static, reduced-form).

\par\smallskip

Empirical evidence of nominal rigidities:

\begin{itemize}[leftmargin=1.7em,itemsep=0.25em,topsep=0.2em]
    \item Taylor (1999): Average price adjustment frequency $=1$ year.
    \item Bils and Klenow (2004): Heterogeneity of frequency across sectors / types of goods.
    \item Nakamura and Steinsson (2008): Adjusting for sales, frequency $\approx 8$--$11$ months.
    \item Bewley (1998): Wages do not fall during recessions (downward nominal rigidity).
\end{itemize}

\par\smallskip

* High welfare / social security countries face stronger downward wage rigidity.

? Working hour rigidity in China.
\[
\text{Nominal rigidity} \Rightarrow \text{price inertia} \Rightarrow
\begin{cases}
\text{Monetary policy changes real balance } M/P,\\
\text{change real interest rate } r=i-\pi.
\end{cases}
\]
\[
\Rightarrow \text{Change in aggregate demand} \Rightarrow \text{Change in employment/output.}
\]
Identify monetary policy shock:
\[
\left\{
\begin{array}{ll}
\text{endogenous movement:} & \text{deliberate response to developments in economy }\Rightarrow\text{co-movement},\\
\text{exogenous movement:} & \text{structural VAR residuals from policy rule.}
\end{array}\right.
\]
\[
\left\{
\begin{array}{@{}l@{\quad}l@{}}
\text{endogenous comovements:}
& \text{high interest rate, high inflation/high output},\\
& i_t=\phi_\pi\pi_t+\phi_y y_t+v_t,\\
& \text{policy response + exogenous component / unexpected shocks};\\[0.35em]
\text{exogenous component:}
& u_t\text{ identification requires: GDP and prices cannot respond}\\
& \text{contemporaneously to a monetary policy shock.}
\end{array}
\right.
\]
\newpage
\section*{Basic NK}
Basic setup:
\begin{enumerate}[label=\circled{\arabic*}]
  \item DSGE structure: Households live infinitely, maximize utility subject to intertemporal budget constraint. Firms with identical technology, subject to random shocks.
  \vspace{0.5em}
  \item AS and AD block.
\end{enumerate}
AD: log-linearized consumption Euler (relates output and price).

\par\smallskip
Let $\hat x_t=\log X_t-\log \bar X$ except $i_t,r_t,\pi_t$.
\[
\tilde y_t=\hat y_t-\hat y_t^{\,n}
\]

$\hat y_t^{\,n}$: natural output under flexible prices / long-run output.
\[
\tilde y_t
=
\E_t\tilde y_{t+1}
-
\frac{1}{\sigma}
\underbrace{\left(\hat i_t-\E_t\hat\pi_{t+1}\right)}_{\hat r_t}
+
\underbrace{v_t}_{\text{MP shock}}.
\tag{1}
\]
AS: New Keynesian Phillips Curve (relates price and output):
\[
\hat\pi_t=\beta \E_t\hat\pi_{t+1}+\kappa\tilde y_t.
\tag{2}
\]
\begin{enumerate}[label=\circled{\arabic*},start=3]
  \item Policy rule:
\end{enumerate}
\[
\hat i_t=f(\Omega_t)+v_t,\qquad \red{\Omega_t:\text{ information set / state vector.}}
\tag{3}
\]
Four markets: goods, labor, bonds, money (cashless).\\
\par\smallskip
Differentiated goods: continuum of varieties $i\in[0,1]$.\\
\par\smallskip
Perfect competition in labor market, no wage rigidities.\\
\par\smallskip
Monopolistic competition on goods market: goods prices set by firms that produce them (profit $\neq 0$).\\
\par\smallskip
Nominal prices are sticky: a fraction $(1-\theta)$ of firms can reset their prices each period (Calvo pricing).\\
\par\smallskip
A continuum of infinitely-lived households.\\
\par\smallskip
A continuum of firms represented by $[0,1]$ goods $\Rightarrow$ \red{1 firm produces 1 differentiated good and sets price.}

\section*{(i) Household's problem}
\[
\max_{C_t,N_t,B_t}\;\E_t\sum_{t=0}^{\infty}\beta^tU(C_t,N_t),
\]
\[
\text{s.t.}\quad \boxit{P_t}C_t+Q_tB_t=B_{t-1}+W_tN_t+T_t.
\]
\blue{$P_t$: aggregate price index of the consumption bundle, $P_t C_t=\int_0^1P_t(i)C_t(i)\dd i$.}
\[
C_t=\left(\int_0^1 C_t(i)^{\frac{\varepsilon-1}{\varepsilon}}\dd i\right)^{\frac{\varepsilon}{\varepsilon-1}},
\]
\red{$C_t$: A composite bundle of goods; continuum of varieties produced by continuum of firms.}
\[
\begin{aligned}
i &:\ \text{variety (differentiated goods) index},\\
\mathrm{d}i &:\ \text{continuum of varieties produced by continuum of firms},\\
\varepsilon &:\ \text{across-variety EoS.}
\end{aligned}
\]
\par\smallskip
\red{Solve household's problem in two steps:}
\par\smallskip
\red{1. Static optimal expenditure allocation (how to allocate across varieties?)}
\[
E=\min_{C_t(i)}\int_0^1P_t(i)C_t(i)\dd i=\text{total expenditure}
\quad \text{s.t.}\quad
C_t=\left(\int_0^1 C_t(i)^{\frac{\varepsilon-1}{\varepsilon}}\dd i\right)^{\frac{\varepsilon}{\varepsilon-1}}.
\]
\[
\mathcal L=\int_0^1P_t(i)C_t(i)\dd i-P_t\left[\left(\int_0^1 C_t(i)^{\frac{\varepsilon-1}{\varepsilon}}\dd i\right)^{\frac{\varepsilon}{\varepsilon-1}}-C_t\right].
\]
\begin{tightalign}
\frac{\partial\mathcal L}{\partial C_t(i)}=0
&\Rightarrow
P_t(i)
-
P_t\frac{\varepsilon}{\varepsilon-1}
\left[
\int_0^1 C_t(i)^{\frac{\varepsilon-1}{\varepsilon}}\dd i
\right]^{\frac{1}{\varepsilon-1}}
\frac{\varepsilon-1}{\varepsilon}
C_t(i)^{-\frac{1}{\varepsilon}}
=0
\\[0.6em]
&\Rightarrow
P_t(i)-P_tC_t^{1/\varepsilon}C_t(i)^{-1/\varepsilon}=0
\\[0.6em]
&\Rightarrow
C_t(i)=
\left(\frac{P_t(i)}{P_t}\right)^{-\varepsilon}C_t.
\end{tightalign}
\begin{tightalign}
C_t(i)^{\frac{\varepsilon-1}{\varepsilon}}
&=\left(\frac{P_t(i)}{P_t}\right)^{1-\varepsilon}C_t^{\frac{\varepsilon-1}{\varepsilon}},\\
\int_0^1C_t(i)^{\frac{\varepsilon-1}{\varepsilon}}\dd i
&=C_t^{\frac{\varepsilon-1}{\varepsilon}}\int_0^1\left(\frac{P_t(i)}{P_t}\right)^{1-\varepsilon}\dd i,\\
C_t^{\frac{\varepsilon-1}{\varepsilon}}
&=C_t^{\frac{\varepsilon-1}{\varepsilon}}\int_0^1\left(\frac{P_t(i)}{P_t}\right)^{1-\varepsilon}\dd i,\\
P_t^{1-\varepsilon}&=\int_0^1P_t(i)^{1-\varepsilon}\dd i,
\qquad
P_t=\left(\int_0^1P_t(i)^{1-\varepsilon}\dd i\right)^{\frac{1}{1-\varepsilon}}.
\end{tightalign}
\blue{Why $\varepsilon$ is the elasticity of substitution across varieties?}
\[
C_t(i)=\left(\frac{P_t(i)}{P_t}\right)^{-\varepsilon}C_t,
\qquad
C_t(j)=\left(\frac{P_t(j)}{P_t}\right)^{-\varepsilon}C_t
\Rightarrow
\frac{C_t(i)}{C_t(j)}=\left(\frac{P_t(i)}{P_t(j)}\right)^{-\varepsilon}.
\]
\[
\log\frac{C_t(i)}{C_t(j)}=-\varepsilon\log\frac{P_t(i)}{P_t(j)},
\qquad
\frac{\dd\log\{C_t(i)/C_t(j)\}}{\dd\log\{P_t(i)/P_t(j)\}}=-\varepsilon.
\]
\[
\sigma_{ij}=-\frac{\dd\log\{C_t(i)/C_t(j)\}}{\dd\log\{P_t(i)/P_t(j)\}}=\varepsilon,
\quad \text{elasticity of substitution between two varieties.}
\]
\par\smallskip
\red{2. Intertemporal problem}
\[
\max_{C_t,N_t,B_t}\E_t\sum_{t=0}^{\infty}\beta^tU(C_t,N_t)
\quad\text{s.t.}\quad P_tC_t+Q_tB_t=B_{t-1}+W_tN_t+T_t.
\]

\[
\mathcal L=\E_t\sum_{t=0}^{\infty}\beta^t\left\{U(C_t,N_t)+\lambda_t(B_{t-1}+W_tN_t+T_t-P_tC_t-Q_tB_t)\right\}.
\]
\begin{tightalign}
\frac{\partial\mathcal L}{\partial C_t}=0&\Rightarrow U_{C,t}=P_t\lambda_t \quad (1),\\
\frac{\partial\mathcal L}{\partial N_t}=0&\Rightarrow U_{N,t}+\lambda_tW_t=0 \quad (2),\\
\frac{\partial\mathcal L}{\partial B_t}=0&\Rightarrow \beta^{t+1}\E_t\lambda_{t+1}-\beta^t\lambda_tQ_t=0
\Rightarrow \E_t\beta\lambda_{t+1}=\lambda_tQ_t \quad (3).
\end{tightalign}
\[(1)\&(2)\Rightarrow -\frac{U_{N,t}}{U_{C,t}}=\frac{W_t}{P_t}\quad\text{Labor--leisure condition (intratemporal).}
\]
\[(1)\&(3)\Rightarrow Q_t=\beta\E_t\left\{\frac{U_{C,t+1}}{U_{C,t}}\frac{P_t}{P_{t+1}}\right\}
=\beta\E_t\left\{\frac{U_{C,t+1}}{U_{C,t}}\frac{1}{\Pi_{t+1}}\right\},
\quad \Pi_{t+1}=\frac{P_{t+1}}{P_t}\text{ inflation}.
\]
\centerline{Bond Euler equation (intertemporal).}
\par\smallskip
Suppose we have a separable utility function
\[
U(C_t,N_t)=\frac{C_t^{1-\sigma}}{1-\sigma}-\frac{N_t^{1+\varphi}}{1+\varphi},
\qquad U_{C,t}=C_t^{-\sigma},\quad U_{N,t}=-N_t^\varphi.
\]
\[
\frac{W_t}{P_t}=N_t^{\varphi}C_t^{\sigma},
\qquad
Q_t=\beta\E_t\left\{\frac{C_{t+1}^{-\sigma}}{C_t^{-\sigma}}\frac{1}{\Pi_{t+1}}\right\}.
\]
Log-linearize:
\[
\hat w_t-\hat p_t=\sigma\hat c_t+\varphi\hat n_t,
\qquad
\hat c_t=\E_t\hat c_{t+1}-\frac{1}{\sigma}(\hat i_t-\hat\Pi_{t+1}).
\]
\[
Q_t=e^{-i_t},
\qquad
\hat\Pi_{t+1}=\log\Pi_{t+1}-\log\bar\Pi=\log\Pi_{t+1}=\pi_{t+1}.
\]
\subsection*{(ii) Firm's problem: two simultaneous problems (isomorphic)}
\red{1. Cost minimization for each producer in the continuum}
\[
\min_{N_t(i)} W_tN_t(i)
\quad\text{s.t.}\quad Y_t(i)=A_tN_t(i)^{1-\alpha}.
\]
\[
\mathcal L=W_tN_t(i)+\mu_t\left[Y_t(i)-A_tN_t(i)^{1-\alpha}\right].
\]
\[
\frac{\partial\mathcal L}{\partial N_t(i)}=W_t-\mu_t(1-\alpha)A_tN_t(i)^{-\alpha}=0
\Rightarrow W_t=\mu_t(1-\alpha)A_tN_t(i)^{-\alpha}.
\]
\[
\mu_t\text{ is the Lagrangian multiplier: shadow value of relaxing the production requirement by one unit.}
\]
\[
\Rightarrow \mu_t \text{ is the nominal marginal cost: } \mu_t=MC_t,
\qquad
\mu_t=MC_t=\frac{W_t}{MPL_t(i)}.
\]

\noindent
$W_t$: 1 labor's cost; \quad
$MPL_t(i)$: 1 labor's output; \quad
$\frac{W_t}{MPL_t(i)}$: cost per unit of output.

\par\smallskip
\par\smallskip
\red{2. Profit maximization by choosing price.}
\[
\max_{P_t(i)}
\;
P_t(i)Y_t(i)-W_tN_t(i)
=
\max_{P_t(i)}
\;
P_t(i)Y_t(i)
-
W_t
\left(
\frac{Y_t(i)}{A_t}
\right)^{\frac{1}{1-\alpha}}
\qquad
\red{\text{production function}}
\]

\[
\frac{\partial Profit}{\partial P_t(i)}
=
Y_t(i)
+
P_t(i)\frac{\partial Y_t(i)}{\partial P_t(i)}
-
W_t
\frac{1}{1-\alpha}
\left(
\frac{Y_t(i)}{A_t}
\right)^{\frac{\alpha}{1-\alpha}}
\frac{\partial Y_t(i)}{\partial P_t(i)}
\frac{1}{A_t}
=0.
\tag{*}
\]
\[
MC_t
=
\frac{\partial}{\partial Y_t(i)}
\left[
W_t
\left(
\frac{Y_t(i)}{A_t}
\right)^{\frac{1}{1-\alpha}}
\right]=W_t
\frac{1}{1-\alpha}
\left(
\frac{Y_t(i)}{A_t}
\right)^{\frac{\alpha}{1-\alpha}}
\frac{1}{A_t}
\]
\[
Y_t(i)
+
\left[
P_t(i)-MC_t
\right]
\frac{\partial Y_t(i)}{\partial P_t(i)}
=0.
\]
Plug in:
\[
Y_t(i)
-
\varepsilon
\left[
P_t(i)-MC_t
\right]
\frac{Y_t(i)}{P_t(i)}
=0.
\]

Divide both sides by \(Y_t(i)\):
\[
1
-
\varepsilon
\frac{P_t(i)-MC_t}{P_t(i)}
=0.
\]

\[
1-\varepsilon+\varepsilon\frac{MC_t}{P_t(i)}
=0.
\]

Therefore,
\[
P_t(i)
=
\frac{\varepsilon}{\varepsilon-1}MC_t.
\]
\par\smallskip
\red{Variety price equals to a constant markup $\mu$ over the marginal cost}

\[
\red{
\text{where }
\mu=\frac{\varepsilon}{\varepsilon-1};
\qquad
\text{monopolistic competition pricing.}
}
\]

\par\smallskip
\blue{When $\varepsilon\to\infty\Rightarrow \mu\to1$, perfect competition, $P_t(i)=MC_t$.}

\blue{In general $\varepsilon$ is finite, $\mu>1$, $P_t(i)>MC_t\Rightarrow$ inefficiency, low employment and output.}
\par\smallskip
\par\smallskip
\red{Cost minimization 1 and profit maximization 2 are dual problems:}
\begin{enumerate}
  \item \red{tells how expensive it is to produce at each output level.}
  \item \red{tells the price-taking that cost schedule and demand as given.}
\end{enumerate}
\red{Nesting the two:}
\[
C(Y_t(i))=\min_{N_t(i)}W_tN_t(i)
\quad\text{s.t.}\quad Y_t(i)=A_tN_t(i)^{1-\alpha}.
\]
Then solve
\[
\max_{P_t(i)}\{P_t(i)Y_t(i)-C(Y_t(i))\}
\quad\text{s.t.}\quad
Y_t(i)=\left(\frac{P_t(i)}{P_t}\right)^{-\varepsilon}C_t.
\]
When $\alpha=0$, marginal cost connects the two problems:
\[
MC_t=\frac{\partial C(Y_t(i))}{\partial Y_t(i)}.
\]
\subsection*{(iii) Sticky-price equilibrium}
Calvo (1983) pricing: Each period a measure $(1-\theta)$ fraction of the continuum of firms reset their prices, the other $\theta$ keeps price unchanged.
\[
\Rightarrow \text{Average duration of a price is }\frac{1}{1-\theta}.
\]
\[
\theta:\text{ index of price stickiness, pinned down by data.}
\]
Thus
\[
P_t=\left[\theta P_{t-1}^{1-\varepsilon}+(1-\theta)(P_t^*)^{1-\varepsilon}\right]^{\frac{1}{1-\varepsilon}}.
\]

Log-linearize:
\begin{tightalign}
P_t^{1-\varepsilon}&=\theta P_{t-1}^{1-\varepsilon}+(1-\theta)(P_t^*)^{1-\varepsilon},\\
\bar P^{1-\varepsilon}(1+\hat p_t)^{1-\varepsilon}
&=\theta\bar P^{1-\varepsilon}(1+\hat p_{t-1})^{1-\varepsilon}+(1-\theta)\bar P^{1-\varepsilon}(1+\hat p_t^*)^{1-\varepsilon},\\
\bar P^{1-\varepsilon}(1-\varepsilon)\hat p_t
&=\theta\bar P^{1-\varepsilon}(1-\varepsilon)\hat p_{t-1}+(1-\theta)\bar P^{1-\varepsilon}(1-\varepsilon)\hat p_t^*,\\
\hat p_t&=\theta\hat p_{t-1}+(1-\theta)\hat p_t^*,\\
\hat p_t-\hat p_{t-1}&=(1-\theta)(\hat p_t^*-
\hat p_{t-1}).
\end{tightalign}
\begin{tightalign}
\hat\pi_t&=\pi_t-\bar\pi=\pi_t=\log P_t-\log P_{t-1}
=\log P_t-\log\bar P-(\log P_{t-1}-\log\bar P)\\
&=\hat p_t-\hat p_{t-1}.
\end{tightalign}
\[
\Rightarrow \hat\pi_t=(1-\theta)(\hat p_t^*-
\hat p_{t-1}).
\]

\par\smallskip
\red{* Positive inflation arises iff firms that adjust prices in $t$ choose to charge above the price level prevailed in $t-1$.}
\[
\red{\theta\uparrow\Rightarrow \text{more sticky firms}\Rightarrow \text{inflation is lower given the price difference.}}
\]

\par\smallskip
\blue{Due to sticky price, firms must set price taking into account the risk that they will not be able to change price in the future.}
\begin{tightalign}
V_t(i)&=P_t(i)Y_t(i)-cost_t(i)+\theta\E_t\{Q_{t,t+1}[P_t(i)Y_{t+1}(i)-cost_{t+1}(i)]\}\\
&\quad+\theta^2\E_t\{Q_{t,t+2}[P_t(i)Y_{t+2}(i)-cost_{t+2}(i)]\}+\cdots.
\end{tightalign}

\par\smallskip
\red{Optimal pricing under sticky price:}
Choose $P_t^*$ to maximize the expected PV of profits while this price remains effective, $P_t^*=P_t^*(i)$.
\[
V_t(i)=\max_{P_t^*}\sum_{k=0}^{\infty}\theta^k\E_t\left[Q_{t,t+k}(P_t^*Y_{t+k|t}-\psi_{t+k}(Y_{t+k|t}))\right]
\quad\purple{\text{conditional on the firm can reset at }t}.
\]
\[
\text{s.t.}\quad Y_{t+k|t}=\left(\frac{P_t^*}{P_{t+k}}\right)^{-\varepsilon}C_{t+k}.
\]
\[
Q_{t,t+k}
=
\beta^k
\left(\frac{C_{t+k}}{C_t}\right)^{-\sigma}
\left(\frac{P_t}{P_{t+k}}\right)
\quad
\text{stochastic discount factor for nominal payoffs.}
\]

\noindent
$\psi_{t+k}(\cdot)$ is the cost function.

\noindent
$Y_{t+k|t}$ is output in period $t+k$ for a firm that last reset price at $t$.

\par\smallskip
\blue{When $\theta=0$, collapse to with only $k=0$ term $\Rightarrow P_t^*=\frac{\varepsilon}{\varepsilon-1}\psi_t$.}\


\par\smallskip
\blue{* $1-\theta$ determines how many firms reset in the cross section.}\\
\blue{$\theta^k$ matters for how long a chosen reset price survives in the time dimension.}

FOC:
\begin{tightalign}
\frac{\partial V_t(i)}{\partial P_t^*}=0
&\Rightarrow
\E_t\left\{\sum_{k=0}^{\infty}\theta^kQ_{t,t+k}\left[Y_{t+k|t}+P_t^*\frac{\partial Y_{t+k|t}}{\partial P_t^*}
-\frac{\partial\psi_{t+k}}{\partial Y_{t+k|t}}\frac{\partial Y_{t+k|t}}{\partial P_t^*}\right]\right\}=0,\\
\frac{\partial Y_{t+k|t}}{\partial P_t^*}
&=-\varepsilon\left(\frac{P_t^*}{P_{t+k}}\right)^{-\varepsilon-1}\frac{C_{t+k}}{P_{t+k}}=-\varepsilon\frac{Y_{t+k|t}}{P_t^*}.
\end{tightalign}
Define
\[
\varphi_{t+k|t}=\frac{\partial\psi_{t+k}}{\partial Y_{t+k|t}}
\quad\text{as nominal marginal cost.}
\]
\begin{tightalign}
&\E_t\left\{\sum_{k=0}^{\infty}\theta^kQ_{t,t+k}\left[Y_{t+k|t}+(P_t^*-\varphi_{t+k|t})(-\varepsilon)\frac{Y_{t+k|t}}{P_t^*}\right]\right\}=0,\\
&\E_t\left\{\sum_{k=0}^{\infty}\theta^kQ_{t,t+k}Y_{t+k|t}\left[1-\varepsilon+\varepsilon\frac{\varphi_{t+k|t}}{P_t^*}\right]\right\}=0.
\end{tightalign}
\[
P_t^*=\underbrace{\frac{\varepsilon}{\varepsilon-1}}_{\red{\text{markup }\mu}}\cdot
\underbrace{\frac{\E_t\sum_{k=0}^{\infty}\theta^kQ_{t,t+k}Y_{t+k|t}\varphi_{t+k|t}}{\E_t\sum_{k=0}^{\infty}\theta^kQ_{t,t+k}Y_{t+k|t}}}_{\red{\text{weighted}}}
\]
\[
\blue{\text{Weighted average of future marginal cost(*).}}
\]
\[
\frac{
\theta^k Q_{t,t+k}Y_{t+k|t}
}{
\E_t\sum_{k=0}^{\infty}\theta^k Q_{t,t+k}Y_{t+k|t}
}
:
\red{\text{weight/how nominal MC of }t+k\text{ matters.}}
\]
When $\theta\to0$, $P_t^*\to\frac{\varepsilon}{\varepsilon-1}\varphi_t$.\\

\par\smallskip
\blue{Weights are larger when:}
\begin{enumerate}[label=\blue{\arabic*)}, leftmargin=2.2em, itemsep=0.15em, topsep=0.15em]
  \item \blue{$\theta\uparrow$: price is more sticky;}
  \item \blue{$Q_{t,t+k}\uparrow$: future profit is more valuable in utility terms;}
  \item \blue{$Y_{t+k|t}\uparrow$: firms expect to sell more in that period.}
\end{enumerate}

\red{Why } 
$Q_{t,t+k}
=
\beta^k
\left(\frac{C_{t+k}}{C_t}\right)^{-\sigma}
\left(\frac{P_t}{P_{t+k}}\right)$?

\[
Q_t
=
\beta
\left[
\left(\frac{C_{t+1}}{C_t}\right)^{-\sigma}
\frac{P_t}{P_{t+1}}
\right]
\quad
\text{one-period nominal discount factor.}
\]

\begin{tightalign}
Q_t
&=
Q_{t,t+1},
\\
Q_{t+1}
&=
\beta
\left[
\left(\frac{C_{t+2}}{C_{t+1}}\right)^{-\sigma}
\frac{P_{t+1}}{P_{t+2}}
\right],
\\
Q_{t,t+2}
&=
Q_{t,t+1}Q_{t+1,t+2}
=
\beta^2
\left(\frac{C_{t+2}}{C_t}\right)^{-\sigma}
\frac{P_t}{P_{t+2}},
\\
Q_{t,t+k}
&=
\beta^k
\left(\frac{C_{t+k}}{C_t}\right)^{-\sigma}
\frac{P_t}{P_{t+k}}.
\end{tightalign}

\noindent
\red{$\beta^k$: impatience.}\\
\red{$\left(\frac{C_{t+k}}{C_t}\right)^{-\sigma}$: marginal utility of future consumption.}\\
\red{$\frac{P_t}{P_{t+k}}$: inflation erosion of nominal payoff; converts nominal payoff into real payoff.}
\[
Q_{t,t+k}=\beta^k\widetilde Q_{t,t+k},
\qquad
\widetilde Q_{t,t+k}=\left(\frac{C_{t+k}}{C_t}\right)^{-\sigma}\left(\frac{P_t}{P_{t+k}}\right).
\]
\[
Y_{t+k|t}=\left(\frac{P_t^*}{P_{t+k}}\right)^{-\varepsilon}C_{t+k}=(P_t^*)^{-\varepsilon}P_{t+k}^{\varepsilon}Y_{t+k},\quad (Y_t=C_t,\forall t).
\]
\[(*)\Rightarrow
P_t^*=\mu\frac{\E_t\sum_{k=0}^{\infty}(\beta\theta)^k\widetilde Q_{t,t+k}P_{t+k}^{\varepsilon}Y_{t+k}\varphi_{t+k|t}}{\E_t\sum_{k=0}^{\infty}(\beta\theta)^k\widetilde Q_{t,t+k}P_{t+k}^{\varepsilon}Y_{t+k}}.
\]
\[
P_t^*\E_t\sum_{k=0}^{\infty}(\beta\theta)^k\widetilde Q_{t,t+k}P_{t+k}^{\varepsilon}Y_{t+k}
=\mu\E_t\sum_{k=0}^{\infty}(\beta\theta)^k\widetilde Q_{t,t+k}P_{t+k}^{\varepsilon}Y_{t+k}\varphi_{t+k|t}.
\]
Let
\[
A_t\equiv\sum_{k=0}^{\infty}(\beta\theta)^k\widetilde Q_{t,t+k}P_{t+k}^{\varepsilon}Y_{t+k},
\qquad
a_{t+k}\equiv(\beta\theta)^k\widetilde Q_{t,t+k}P_{t+k}^{\varepsilon}Y_{t+k}.
\]
\[
\widehat{LHS}=\hat p_t^*+\E_t\hat A_t.
\]
\begin{tightalign}
\hat A_t&=\sum_{k=0}^{\infty}\frac{\bar a_{t+k}}{\bar A}\hat a_{t+k},\\
\frac{\bar a_{t+k}}{\bar A}
&=\frac{(\beta\theta)^k\bar{\widetilde Q}\bar P^{\varepsilon}\bar Y}{\sum_{k=0}^{\infty}(\beta\theta)^k\bar{\widetilde Q}\bar P^{\varepsilon}\bar Y}
=\frac{(\beta\theta)^k\bar{\widetilde Q}\bar P^{\varepsilon}\bar Y}{\frac{1}{1-\beta\theta}\bar{\widetilde Q}\bar P^{\varepsilon}\bar Y}
=(1-\beta\theta)(\beta\theta)^k,\\
\hat a_{t+k}&=\hat{\widetilde q}_{t,t+k}+\varepsilon\hat p_{t+k}+\hat y_{t+k}.
\end{tightalign}

\par\smallskip
\[
\widehat{LHS}=\hat p_t^*+(1-\beta\theta)\E_t\sum_{k=0}^{\infty}(\beta\theta)^k(\hat{\widetilde q}_{t,t+k}+\varepsilon\hat p_{t+k}+\hat y_{t+k}).
\]
\[
B_t\equiv\E_t\sum_{k=0}^{\infty}(\beta\theta)^k\widetilde Q_{t,t+k}P_{t+k}^{\varepsilon}Y_{t+k}\varphi_{t+k|t},
\qquad
\widehat{RHS}=\E_t\hat B_t.
\]
\begin{tightalign}
\hat B_t&=\sum_{k=0}^{\infty}\frac{\bar b_{t+k}}{\bar B}\hat b_{t+k},
\qquad
\frac{\bar b_{t+k}}{\bar B}=(1-\beta\theta)(\beta\theta)^k,\\
\hat b_{t+k}&=\hat{\widetilde q}_{t,t+k}+\varepsilon\hat p_{t+k}+\hat y_{t+k}+\hat\varphi_{t+k|t}.
\end{tightalign}
\[
\widehat{RHS}=(1-\beta\theta)\E_t\sum_{k=0}^{\infty}(\beta\theta)^k(\hat{\widetilde q}_{t,t+k}+\varepsilon\hat p_{t+k}+\hat y_{t+k}+\hat\varphi_{t+k|t}).
\]
Combine $\widehat{LHS}=\widehat{RHS}$:
\begin{tightalign}
\hat p_t^*&=(1-\beta\theta)\E_t\sum_{k=0}^{\infty}(\beta\theta)^k\hat\varphi_{t+k|t}\\
&=(1-\beta\theta)\E_t\sum_{k=0}^{\infty}(\beta\theta)^k(\widehat{mc}^{\ r}_{t+k|t}+\hat p_{t+k}).
\end{tightalign}
\red{nominal MC $=$ real MC $\times$ price}
\[
\varphi_{t+k|t}=MC^r_{t+k|t}\cdot P_{t+k}
\Rightarrow \hat\psi_{t+k|t}=\widehat{mc}^{\ r}_{t+k|t}+\hat p_{t+k}.
\]

\par\smallskip
\blue{1. Weights $\sum_{k=0}^{\infty}(1-\beta\theta)(\beta\theta)^k=1$: weighted average over expected future nominal marginal cost.}

\par\smallskip
\blue{2. Aggregate price level denotes willingness to maintain (in expectation) the relative price unchanged.}
\[
\text{Suppose } \widehat{mc}^{\,r}_{t+k|t}=0 \text{ for all } k,
\quad
\text{real marginal cost is constant.}
\]

\[
\hat p_t^*
=
(1-\beta\theta)
\E_t
\sum_{k=0}^{\infty}
(\beta\theta)^k
\hat p_{t+k}.
\]
\blue{* Reset price rises when expected future aggregate prices rise.}\
\red{Front-load future inflation to today's reset price. Expected inflation $\Rightarrow \hat p_t^*\uparrow$.}

\par\smallskip
\blue{3. Real MC $\widehat{mc}^{\ r}_{t+k|t}$ denotes desire to change the expected relative price to avoid any gap that may emerge between expected and desired markup. Monopolistic competitive firm's desired markup rule: $P=\mu MC$.}
\[
\text{Thus, firm wants its real price }\frac{P_t^*}{P_{t+k}}\text{ to be high enough relative to real MC.}
\]
\[
\hat p_t^* -\hat p_{t+k}=\widehat{mc}^{\ r}_{t+k|t}+\text{desired markup}
\Rightarrow \widehat{mc}^{\ r}_{t+k|t}\uparrow\Rightarrow\hat p_t^* -\hat p_{t+k}\uparrow.
\]
\blue{$\Rightarrow$ Expected future real marginal cost is higher: the firm wants a higher relative price.}

\newpage
\section*{(iv) Equilibrium}
\textbf{1. Goods market clearing:}
\[
Y_t(i)=C_t(i),\quad \forall i\in[0,1].
\]
Define
\[
Y_t=\left(\int_0^1Y_t(i)^{\frac{\varepsilon-1}{\varepsilon}}\dd i\right)^{\frac{\varepsilon}{\varepsilon-1}}.
\]
Get $Y_t=C_t$. Log-linearize: $\hat y_t=\hat c_t$. Since
\[
\hat c_t=\E_t\hat c_{t+1}-\frac{1}{\sigma}(\hat i_t-\E_t\hat\pi_{t+1}),
\]
\[
\Rightarrow \hat y_t=\E_t\hat y_{t+1}-\frac{1}{\sigma}(\hat i_t-\E_t\hat\pi_{t+1}).
\]
\textbf{2. Labor market clearing:}
\[
N_t=\int_0^1N_t(i)\dd i.
\]
\begin{tightalign}
Y_t(i)=A_tN_t(i)^{1-\alpha}
&\Rightarrow N_t=\int_0^1\left(\frac{Y_t(i)}{A_t}\right)^{\frac{1}{1-\alpha}}\dd i,\\
Y_t(i)=\left(\frac{P_t(i)}{P_t}\right)^{-\varepsilon}Y_t
&\Rightarrow N_t=\left(\frac{Y_t}{A_t}\right)^{\frac{1}{1-\alpha}}
\int_0^1\left(\frac{P_t(i)}{P_t}\right)^{-\frac{\varepsilon}{1-\alpha}}\dd i,\\
&\Rightarrow N_t^{1-\alpha}=\frac{Y_t}{A_t}\left[\int_0^1\left(\frac{P_t(i)}{P_t}\right)^{-\frac{\varepsilon}{1-\alpha}}\dd i\right]^{1-\alpha}.
\end{tightalign}
\section*{Price dispersion}

Let
\[
D_t
=
\left[
\int_0^1
\left(\frac{P_t(i)}{P_t}\right)^{-\frac{\varepsilon}{1-\alpha}}
\dd i
\right]^{1-\alpha}
=
\exp(d_t).
\]

\[
Y_t
=
A_tN_t^{1-\alpha}D_t^{-1},
\qquad
\blue{D_t\text{ measures how uneven prices are across firms.}}
\]

\[
\blue{\text{When all firms charge the same price:}}\quad
\frac{P_t(i)}{P_t}=1,\ \forall i
\Rightarrow
D_t=1
\Rightarrow
d_t=0.
\]

\par\smallskip
\noindent
\blue{Reduced to $Y_t=A_tN_t^{1-\alpha}$: frictionless benchmark under flexible prices, no price distortion.}


\par\smallskip
\blue{But under Calvo pricing, some firms use old price, some firms use new price $\Rightarrow D_t>1$.}
\[
Y_t(i)=\left(\frac{P_t(i)}{P_t}\right)^{-\varepsilon}Y_t.
\]

\noindent
\blue{Firms with low relative prices sell too much; firms with high relative prices sell too little.}
\[
d_t=(1-\alpha)\log\left[\int_0^1\left(\frac{P_t(i)}{P_t}\right)^{-\frac{\varepsilon}{1-\alpha}}\dd i\right].
\]
Define $p_t(i)=\log P_t(i)$, $p_t=\log P_t$, $P_t(i)/P_t=e^{p_t(i)-p_t}$. Let $\chi_i=p_t(i)-p_t$.
\[
\Rightarrow d_t=(1-\alpha)\log\left[\int_0^1\exp\left(-\frac{\varepsilon}{1-\alpha}\chi_i\right)\dd i\right].
\]

\par\smallskip
Second-order Taylor expansion around zero-inflation / zero-dispersion steady state.\\
\red{Why zero inflation $\bar\Pi=1\Leftrightarrow$ zero dispersion $P_t(i)/P_t=1\Leftrightarrow \bar d=0$?}
\[
P_t=\left[\theta P_{t-1}^{1-\varepsilon}+(1-\theta)(P_t^*)^{1-\varepsilon}\right]^{\frac{1}{1-\varepsilon}},
\quad \red{\text{to have }P_t=P_{t-1}=\bar P,\text{ need }P_t^*=\bar P.}
\]
\red{The firms that reset price set the same price as others.}
\[
\text{Since } 
e^z
=
e^0+\frac{(e^z)|_{z=0}}{1!}(z-0)
+\frac{(e^z)|_{z=0}}{2!}(z-0)^2
+o(z^2).
\]

\[
e^z\approx 1+z+\frac{z^2}{2}
\quad \text{for small } z.
\]

\[
e^{-\frac{\varepsilon}{1-\alpha}\chi_i}
\approx
1-\frac{\varepsilon}{1-\alpha}\chi_i
+
\frac{1}{2}
\left(\frac{\varepsilon}{1-\alpha}\right)^2
\chi_i^2.
\]

\begin{tightalign}
\int_0^1 e^{-\frac{\varepsilon}{1-\alpha}\chi_i}\dd i
&\approx
1-\frac{\varepsilon}{1-\alpha}\int_0^1\chi_i\dd i
+
\frac{1}{2}
\left(\frac{\varepsilon}{1-\alpha}\right)^2
\int_0^1\chi_i^2\dd i,\\
\chi_i=p_t(i)-p_t
&\Rightarrow
\int_0^1\chi_i\dd i
=
\int_0^1p_t(i)\dd i-\int_0^1p_t\dd i
=
p_t-p_t=0,\\
\int_0^1e^{-\frac{\varepsilon}{1-\alpha}\chi_i}\dd i
&\approx
1+
\frac{1}{2}
\left(\frac{\varepsilon}{1-\alpha}\right)^2
\int_0^1\chi_i^2\dd i.
\end{tightalign}

\[
\int_0^1\chi_i^2\dd i
\approx
\operatorname{Var}_i\{p_t(i)-p_t\}
\quad
\text{cross-sectional second moment.}
\]
Take logs:
\begin{tightalign}
d_t&=(1-\alpha)\log\left[1+\frac{1}{2}\left(\frac{\varepsilon}{1-\alpha}\right)^2Var_i\{p_t(i)-p_t\}\right],\\
\text{since }\log(1+u)&\approx u,\\
d_t&\approx (1-\alpha)\frac{1}{2}\left(\frac{\varepsilon}{1-\alpha}\right)^2Var_i\{p_t(i)-p_t\},\\
d_t&\approx \frac{\varepsilon^2}{2(1-\alpha)}Var_i\{p_t(i)-p_t\},\\
d_t&\propto Var_i\{p_t(i)\}.
\end{tightalign}
\par\smallskip
Thus, in the neighborhood of $\hat\pi_t=0$ or $\bar\Pi=1$, $d_t$ is $0$ up to first-order approximation:
\[
\hat d_t\approx0.
\]
Thus
\[
Y_t=A_tN_t^{1-\alpha}D_t^{-1}\quad\text{log linearization}
\]
\[
\hat y_t=\hat a_t+(1-\alpha)\hat n_t-\hat d_t
\Rightarrow
\hat y_t=\hat a_t+(1-\alpha)\hat n_t,
\qquad
\hat n_t=\frac{1}{1-\alpha}(\hat y_t-\hat a_t).
\]
\newpage
\section*{Average real marginal cost}
\red{Average firm:}
\[
MC_t^r=\frac{W_t}{P_t}\frac{1}{MPN_t}=\frac{W_t}{P_t}\frac{1}{(1-\alpha)A_tN_t^{-\alpha}}.
\]
\begin{tightalign}
\widehat{mc}_t^r
&=\hat w_t-\hat p_t-\widehat{mpn}_t
=\hat w_t-\hat p_t-(\hat a_t-\alpha\hat n_t)\\
&=\hat w_t-\hat p_t-\left(\hat a_t-\frac{\alpha}{1-\alpha}\hat y_t+\frac{\alpha}{1-\alpha}\hat a_t\right)\\
&=\hat w_t-\hat p_t-\frac{1}{1-\alpha}(\hat a_t-\alpha\hat y_t).
\end{tightalign}
\red{For firms that set price at $t$ and remains at $t+k$:}
\[
\begin{aligned}
\widehat{mc}^{\,r}_{t+k|t}
&=
\hat w_{t+k}
-
\hat p_{t+k}
-
\widehat{mpn}_{t+k|t}
\\[0.15em]
&=
\hat w_{t+k}
-
\hat p_{t+k}
-
\frac{1}{1-\alpha}
\left(
\hat a_{t+k}
-
\alpha \hat y_{t+k|t}
\right)
\\[0.15em]
&=
\hat w_{t+k}
-
\hat p_{t+k}
-
\frac{1}{1-\alpha}
\left(
\hat a_{t+k}
-
\alpha \hat y_{t+k}
+
\alpha \hat y_{t+k}
-
\alpha \hat y_{t+k|t}
\right)
\\[0.15em]
&=
\widehat{mc}^{\,r}_{t+k}
-
\frac{\alpha}{1-\alpha}
\left(
\hat y_{t+k}
-
\hat y_{t+k|t}
\right).
\end{aligned}
\]
Since
\[
Y_{t+k|t}=\left(\frac{P_t^*}{P_{t+k}}\right)^{-\varepsilon}Y_{t+k}
\Rightarrow \hat y_{t+k}-\hat y_{t+k|t}=\varepsilon(\hat p_t^*-
\hat p_{t+k}).
\]
\[
\Rightarrow\widehat{mc}_{t+k|t}^r=\widehat{mc}_{t+k}^r-\frac{\alpha\varepsilon}{1-\alpha}(\hat p_t^*-
\hat p_{t+k}).
\]
\par\smallskip
\blue{With CRS: $\alpha=0\Rightarrow\widehat{mc}_{t+k|t}^r=\widehat{mc}_{t+k}^r$. The real marginal cost is common across firms regardless whether they can reset price at $t$.}

Since
\[
\widehat{mc}^{\,r}_{t+k|t}
=
\widehat{mc}^{\,r}_{t+k}
-
\frac{\alpha\varepsilon}{1-\alpha}
(\hat p_t^*-\hat p_{t+k}),
\]
we have
\[
\hat p_t^*
=
(1-\beta\theta)E_t
\sum_{k=0}^{\infty}
(\beta\theta)^k
\left[
\Theta \widehat{mc}^{\,r}_{t+k}
+
\hat p_{t+k}
\right],
\qquad
\Theta=
\frac{1-\alpha}{1-\alpha+\alpha\varepsilon}\leq 1.
\]
Derivation:
\begin{tightalign}
\hat p_t^*&=(1-\beta\theta)\E_t\left\{\sum_{k=0}^{\infty}(\beta\theta)^k\left[\widehat{mc}_{t+k}^r-\frac{\alpha\varepsilon}{1-\alpha}(\hat p_t^*-
\hat p_{t+k})+\hat p_{t+k}\right]\right\},\\
\hat p_t^*&=-\frac{\alpha\varepsilon}{1-\alpha}\hat p_t^*+(1-\beta\theta)\E_t\left\{\sum_{k=0}^{\infty}(\beta\theta)^k\left[\widehat{mc}_{t+k}^r+\frac{1-\alpha+\alpha\varepsilon}{1-\alpha}\hat p_{t+k}\right]\right\},\\
\frac{1-\alpha+\alpha\varepsilon}{1-\alpha}\hat p_t^*&=(1-\beta\theta)\E_t\left\{\sum_{k=0}^{\infty}(\beta\theta)^k\left[\widehat{mc}_{t+k}^r+\frac{1-\alpha+\alpha\varepsilon}{1-\alpha}\hat p_{t+k}\right]\right\},\\
\hat p_t^*&=(1-\beta\theta)\E_t\left\{\sum_{k=0}^{\infty}(\beta\theta)^k[\Theta\widehat{mc}_{t+k}^r+\hat p_{t+k}]
\right\}.
\end{tightalign}
\begin{tightalign}
\hat p_t^*&=(1-\beta\theta)(\Theta\widehat{mc}_t^r+\hat p_t)+\beta\theta\E_t\hat p_{t+1}^* \quad (*),\\
\beta\theta\hat p_{t+1}^*&=\beta\theta(1-\beta\theta)\E_t\left\{\sum_{k=0}^{\infty}(\beta\theta)^k[\Theta\widehat{mc}_{t+1+k}^r+\hat p_{t+1+k}]
\right\}\\
&=(1-\beta\theta)\E_t\left\{\sum_{k=0}^{\infty}(\beta\theta)^{k+1}[\Theta\widehat{mc}_{t+(k+1)}^r+\hat p_{t+(k+1)}]
\right\}.
\end{tightalign}
Let $k'=k+1$, $\forall k$:
\[
\beta\theta\hat p_{t+1}^*=(1-\beta\theta)\E_t\left\{\sum_{k'=1}^{\infty}(\beta\theta)^{k'}[\Theta\widehat{mc}_{t+k'}^r+\hat p_{t+k'}]\right\}.
\]
\blue{$k'$ is any index; we can use $k$.}\
\blue{(*) just take out the term of $k=0$.}
Combine (*) with log-linearized price index:
\[
\hat p_t=\theta\hat p_{t-1}+(1-\theta)\hat p_t^*.
\]
\begin{tightalign}
\hat p_t&=\theta\hat p_{t-1}+(1-\theta)(1-\beta\theta)(\Theta\widehat{mc}_t^r+
\hat p_t)+(1-\theta)\beta\theta\E_t\hat p_{t+1}^*,\\
\text{since }\hat p_{t+1}^*&=\frac{\E_t\hat p_{t+1}-\theta\hat p_t}{1-\theta},\\
\hat p_t&=\theta\hat p_{t-1}+(1-\theta)(1-\beta\theta)(\Theta\widehat{mc}_t^r+
\hat p_t)+(1-\theta)\beta\theta\frac{\E_t\hat p_{t+1}-\theta\hat p_t}{1-\theta}.
\end{tightalign}

\begin{tightalign}
\hat p_t&=\theta\hat p_{t-1}+(1-\theta)(1-\beta\theta)\Theta\widehat{mc}_t^r+(1-\theta-\beta\theta+\beta\theta^2)\hat p_t+\beta\theta\E_t\hat p_{t+1}-\beta\theta^2\hat p_t,\\
\hat p_t-\theta\hat p_{t-1}-(1-\theta)\hat p_t&=(1-\theta)(1-\beta\theta)\Theta\widehat{mc}_t^r+\beta\theta(\E_t\hat p_{t+1}-\hat p_t),\\
\theta(\hat p_t-\hat p_{t-1})&=(1-\theta)(1-\beta\theta)\Theta\widehat{mc}_t^r+\beta\theta\E_t\hat\pi_{t+1},\\
\hat\pi_t&=\underbrace{\frac{(1-\theta)(1-\beta\theta)}{\theta}\Theta}_{\lambda}\widehat{mc}_t^r+\beta\E_t\hat\pi_{t+1}.
\end{tightalign}
\[
\Rightarrow \hat\pi_t=\beta\E_t\hat\pi_{t+1}+\lambda\widehat{mc}_t^r\quad \text{iterate forward}
\]
\[
\Rightarrow \hat\pi_t=\lambda\sum_{k=0}^{\infty}\beta^k\E_t\widehat{mc}_{t+k}^r.
\]
\par\smallskip
\noindent
\blue{$\lambda$ is the slope of NKPC w.r.t. real marginal cost.}\\
\blue{$\theta\uparrow$ sticky price $\Rightarrow \lambda\downarrow \Rightarrow$ flatter Phillips curve.}\\
\blue{$\Theta\downarrow$ stronger decreasing returns $\Rightarrow \lambda\downarrow$.}
\par\smallskip
\noindent
\begin{minipage}{0.96\linewidth}
\par\smallskip
\red{Inflation results from aggregate consequences of purposeful price-setting decisions by firms,
which adjust their prices in light of current and anticipated cost conditions.}
\end{minipage}
\par\smallskip
\par\smallskip
\section*{Relate marginal cost with output gap}
\begin{tightalign}
\widehat{mc}_t^r
&=\hat w_t-\hat p_t-\frac{1}{1-\alpha}(\hat a_t-\alpha\hat y_t),\\
\hat w_t-
\hat p_t&=\sigma\hat c_t+\varphi\hat n_t
\Rightarrow \hat w_t-
\hat p_t=\sigma\hat y_t+\varphi\hat n_t,\\
\Rightarrow \widehat{mc}_t^r&=\frac{\varphi+\alpha+\sigma-\alpha\sigma}{1-\alpha}\hat y_t-\frac{\varphi+1}{1-\alpha}\hat a_t.
\end{tightalign}
\par\smallskip
\red{Real marginal cost rises with output, falls with productivity.}
\par\smallskip
\begin{enumerate}
  \item \red{Labor becomes harder to expand due to household requires more compensation.}
  \item \red{Decreasing return to scale: labor becomes less productive at the margin.}
\end{enumerate}
\par\smallskip
Under flexible price:
\[
P_t(i)=\frac{\varepsilon}{\varepsilon-1}\frac{W_t}{A_tF_N,t(N_t(i))}=\mu MC_t=P_t,
\]
\[
MC_t^r=\frac{MC_t}{P_t}=\frac{1}{\mu}\Rightarrow \widehat{mc}_t^r=0,
\quad \text{real marginal cost has no variation.}
\]
\par\smallskip
Define the output under flexible price as natural output (long-run output) $Y_t^n$.
\[
\Rightarrow \text{log deviation }\hat y_t^n\Rightarrow
\hat y_t^n=\frac{\varphi+1}{\varphi+\alpha+\sigma-\alpha\sigma}\hat a_t.
\]
\par\smallskip
\red{Define output gap as } $\tilde y_t=\hat y_t-\hat y_t^n$.
\begin{tightalign}
\text{Then }\widehat{mc}_t^r
&=\frac{\varphi+\alpha+\sigma-\alpha\sigma}{1-\alpha}(\hat y_t^n+\tilde y_t)-\frac{\varphi+1}{1-\alpha}\hat a_t\\
&=\frac{\varphi+\alpha+\sigma-\alpha\sigma}{1-\alpha}\frac{\varphi+1}{\varphi+\alpha+\sigma-\alpha\sigma}\hat a_t+\frac{\varphi+\alpha+\sigma-\alpha\sigma}{1-\alpha}\tilde y_t-\frac{\varphi+1}{1-\alpha}\hat a_t\\
&=\frac{\varphi+\alpha+\sigma(1-\alpha)}{1-\alpha}\tilde y_t
=\left(\sigma+\frac{\varphi+\alpha}{1-\alpha}\right)\tilde y_t.
\end{tightalign}
NKPC:
\[
\hat\pi_t=\beta\E_t\hat\pi_{t+1}+\lambda\widehat{mc}_t^r\quad(\hat\pi_{t+1}=\pi_{t+1})
\]
\[
\Rightarrow \hat\pi_t=\beta\E_t\hat\pi_{t+1}+\kappa\tilde y_t,
\quad \kappa=\lambda\left(\sigma+\frac{\varphi+\alpha}{1-\alpha}\right),
\quad \lambda=\frac{(1-\theta)(1-\beta\theta)}{\theta}\Theta,
\quad \Theta=\frac{1-\alpha}{1-\alpha+\alpha\varepsilon}.
\]
\par\smallskip
\blue{When demand pushes output above natural output, firms face tighter resource constraints, and the marginal cost rises.}
\begin{enumerate}
  \item \blue{$\sigma$: demand / consumption smoothing channel.}
  \[
  Y_t>Y_t^n\Rightarrow C_t>C_t^n\Rightarrow MU\text{ of leisure}\uparrow\Rightarrow\text{tends to work less, but to keep }N\Rightarrow\frac{W}{P}\uparrow\Rightarrow MC^r\uparrow.
  \]
 \[
\begin{aligned}
&\sigma\uparrow
\Rightarrow \text{less willing to substitute consumption intertemporally} \\
&\Rightarrow \text{stronger intertemporal } C,N \text{ substitution} \\
&\Rightarrow \text{stronger wage pressure}
\Rightarrow \text{stronger MC pressure.}
\end{aligned}
\]
  \item \blue{$\varphi$: labor market channel.}
  \[
  \varphi\uparrow\Rightarrow\text{households dislike work}\Rightarrow W\text{ increases by more}\Rightarrow MC^r\text{ increases by more.}
  \]
  \item \blue{$\alpha$: diminishing-returns-to-scale channel.}
  \[
  \alpha\uparrow\Rightarrow \text{to increase output by 1 unit}\Rightarrow \alpha\uparrow\text{ units increase in labor}\Rightarrow W/P\uparrow\Rightarrow MC^r\uparrow.
  \]
\end{enumerate}
\newpage
\section*{Dynamic IS curve}
\begin{tightalign}
\hat y_t&=\E_t\{\hat y_{t+1}\}-\frac{1}{\sigma}\left(\hat i_t-\E_t\{\hat\pi_{t+1}\}\right),\\
\hat y_t^n+\tilde y_t&=\E_t\{\hat y_{t+1}^n+\tilde y_{t+1}\}-\frac{1}{\sigma}\left(\hat i_t-\E_t\{\hat\pi_{t+1}\}\right),\\
\tilde y_t&=\E_t\{\hat y_{t+1}^n-\hat y_t^n\}+\E_t\{\tilde y_{t+1}\}-\frac{1}{\sigma}(\hat i_t-\E_t\{\hat\pi_{t+1}\}).
\end{tightalign}
Define natural interest rate as:
\[
\hat r_t^n=\sigma\E_t\{\hat y_{t+1}^n-\hat y_t^n\}
=\frac{\sigma(1+\varphi)}{\varphi+\alpha+\sigma(1-\alpha)}\E_t\{\Delta\hat a_{t+1}\}
=\sigma\psi_{ya}^n\E_t\{\Delta\hat a_{t+1}\}.
\]
\[
\Rightarrow \tilde y_t=-\frac{1}{\sigma}(\hat i_t-\E_t\{\hat\pi_{t+1}\}-\hat r_t^n)+\E_t\{\tilde y_{t+1}\}.
\]
Define
\[
\hat r_t=\hat i_t-\E_t\{\hat\pi_{t+1}\}.
\]
\[
\Rightarrow \tilde y_t=-\frac{1}{\sigma}(\hat r_t-\hat r_t^n)+\E_t\{\tilde y_{t+1}\}.
\]
Assume $\lim_{T\to\infty}\E_t\{\tilde y_{t+T}\}=0$ and iterate forward:
\[
\Rightarrow \tilde y_t=-\frac{1}{\sigma}\E_t\left\{\sum_{k=0}^{\infty}(\hat r_{t+k}-\hat r_{t+k}^n)\right\}.
\]
\blue{Output gap is proportional to the sum of current and anticipated deviations between the real interest rate and natural real interest rate.}
\par\smallskip
NKPC:
\[
\hat\pi_t=\beta\E_t\{\hat\pi_{t+1}\}+\kappa\tilde y_t\quad\text{iterate forward}
\]
\[
\Rightarrow\hat\pi_t=\kappa\sum_{k=0}^{\infty}\beta^k\E_t\{\tilde y_{t+k}\},
\quad \text{where }\kappa=\lambda\left(\sigma+\frac{\varphi+\alpha}{1-\alpha}\right).
\]
\[
\red{\text{Inflation today is the discounted sum of the expected future output gap.}}
\]
\par\smallskip
Basic New Keynesian non-policy block:
\[
\left\{
\begin{array}{ll}
\hat\pi_t=\kappa\sum_{k=0}^{\infty}\beta^k\E_t\{\tilde y_{t+k}\}, & \text{NKPC: determines inflation given output gap path},\\[3pt]
\tilde y_t=-\frac{1}{\sigma}\E_t\left\{\sum_{k=0}^{\infty}(\hat r_{t+k}-\hat r_{t+k}^n)\right\}, & \text{DIS: determines output gap given real interest rate gap path.}
\end{array}
\right.
\]

\par\smallskip
To close, we need more equations to pin down how nominal interest rate $i$ / $\hat i_t$ evolves (monetary policy).
\[
\left\{
\begin{aligned}
\tilde y_t
&=
-\frac{1}{\sigma}
\left(
\hat i_t-\mathbb E_t\{\hat\pi_{t+1}\}-\hat r_t^n
\right)
+\mathbb E_t\{\tilde y_{t+1}\}
\qquad \dashrightarrow \quad \text{DIS (Demand)}
\\[0.6em]
\hat\pi_t
&=
\beta \mathbb E_t\{\hat\pi_{t+1}\}
+\kappa \tilde y_t
\qquad \dashrightarrow \quad \text{NKPC (Supply)}
\end{aligned}
\right.
\]
\newpage
\section*{Dynamic monetary policy rule}
\[
\hat i_t=\phi_\pi\hat\pi_t+\phi_y\tilde y_t+v_t,
\quad v_t\text{ is exogenous with zero mean.}
\]
\[
\phi_\pi>0,\quad \phi_y>0,
\quad \text{assume S.S. }\bar\Pi=0.
\]
DIS:
\[
\tilde y_t=-\frac{1}{\sigma}[\phi_\pi\hat\pi_t+\phi_y\tilde y_t-
\E_t\{\hat\pi_{t+1}\}-(\hat r_t^n-v_t)]+\E_t\{\tilde y_{t+1}\}\quad (***)
\]
\[
\Rightarrow \left(1+\frac{\phi_y}{\sigma}\right)\tilde y_t+\frac{\phi_\pi}{\sigma}\hat\pi_t
=\E_t\{\tilde y_{t+1}\}+\frac{1}{\sigma}\E_t\{\hat\pi_{t+1}\}+\frac{1}{\sigma}(\hat r_t^n-v_t).
\]
\[
-\kappa\tilde y_t+\hat\pi_t=\beta\E_t\{\hat\pi_{t+1}\}.
\]
\par\smallskip
Write in matrix form in system of $[\tilde y_t,\hat\pi_t]'$:
\[
\underbrace{\begin{bmatrix}1+\phi_y/\sigma&\phi_\pi/\sigma\\-\kappa&1\end{bmatrix}}_{M}
\begin{bmatrix}\tilde y_t\\\hat\pi_t\end{bmatrix}
=\begin{bmatrix}1&1/\sigma\\0&\beta\end{bmatrix}
\begin{bmatrix}\E_t\{\tilde y_{t+1}\}\\\E_t\{\hat\pi_{t+1}\}\end{bmatrix}
+\begin{bmatrix}1/\sigma\\0\end{bmatrix}(\hat r_t^n-v_t).
\]
\[
\begin{bmatrix}\tilde y_t\\\hat\pi_t\end{bmatrix}
=M^{-1}\begin{bmatrix}1&1/\sigma\\0&\beta\end{bmatrix}
\begin{bmatrix}\E_t\{\tilde y_{t+1}\}\\\E_t\{\hat\pi_{t+1}\}\end{bmatrix}
+M^{-1}\begin{bmatrix}1/\sigma\\0\end{bmatrix}(\hat r_t^n-v_t).
\]
\[
M^{-1}=\frac{1}{|M|}\begin{bmatrix}1&-\phi_\pi/\sigma\\ \kappa&1+\phi_y/\sigma\end{bmatrix}
=\Omega\begin{bmatrix}\sigma&-\phi_\pi\\ \kappa\sigma&\sigma+\phi_y\end{bmatrix},
\quad \Omega=\frac{1}{\sigma+\phi_y+\kappa\phi_\pi}.
\]
\[
A_T=\Omega\begin{bmatrix}\sigma&-\phi_\pi\\\kappa\sigma&\sigma+\phi_y\end{bmatrix}
\begin{bmatrix}1&1/\sigma\\0&\beta\end{bmatrix}
=\begin{bmatrix}\sigma&1-\beta\phi_\pi\\ \kappa\sigma&\kappa+\beta(\sigma+\phi_y)\end{bmatrix}\Omega,
\]
\[
B_T=\Omega\begin{bmatrix}\sigma&-\phi_\pi\\\kappa\sigma&\sigma+\phi_y\end{bmatrix}
\begin{bmatrix}1/\sigma\\0\end{bmatrix}
=\begin{bmatrix}1\\\kappa\end{bmatrix}\Omega.
\]

\[
\Rightarrow
\begin{bmatrix}\tilde y_t\\\hat\pi_t\end{bmatrix}
=A_T\begin{bmatrix}\E_t\{\tilde y_{t+1}\}\\\E_t\{\hat\pi_{t+1}\}\end{bmatrix}
+B_T(\hat r_t^n-v_t).
\]
\par\smallskip
\[
\red{\text{Both }\tilde y_t\text{ and }\hat\pi_t\text{ are not predetermined }\xRightarrow{BK}\text{ both eigenvalues of }A_T\text{ are within the unit circle.}}
\]
For a $2\times2$ matrix $A_T$, both eigenvalues lie inside unit circle iff Jury/Schur condition holds:
\[
\left\{
\begin{array}{l}
1-\det(A_T)>0\quad (D<1),\\
1-\operatorname{tr}(A_T)+\det(A_T)>0\quad (p(1)>0),\\
1+\operatorname{tr}(A_T)+\det(A_T)>0\quad (p(-1)>0).
\end{array}
\right.
\]
\begin{tightalign}
\det(A_T)&=\Omega^2[\sigma(\kappa+\beta(\sigma+\phi_y))-(1-\beta\phi_\pi)\kappa\sigma]
=\Omega^2\sigma\beta(\sigma+\phi_y+\kappa\phi_\pi)
=\frac{\beta\sigma}{\sigma+\phi_y+\kappa\phi_\pi},\\
\operatorname{tr}(A_T)&=\Omega[\sigma+\kappa+\beta(\sigma+\phi_y)]
=\frac{\sigma+\kappa+\beta(\sigma+\phi_y)}{\sigma+\phi_y+\kappa\phi_\pi}.
\end{tightalign}
\[
1^\circ\quad
1-
\det(A_T)=1-\frac{\beta\sigma}{\sigma+\phi_y+\kappa\phi_\pi}
=\frac{(1-\beta)\sigma+\phi_y+\kappa\phi_\pi}{\sigma+\phi_y+\kappa\phi_\pi}>0\quad \blue{\text{satisfy}},
\]
$\beta<1,\sigma>0,\phi_y>0,\kappa>0,\phi_\pi>0.$
\begin{tightalign}
2^\circ\quad 1-\operatorname{tr}(A_T)+\det(A_T)
&=1-\frac{\sigma+\kappa+\beta(\sigma+\phi_y)}{\sigma+\phi_y+\kappa\phi_\pi}+\frac{\beta\sigma}{\sigma+\phi_y+\kappa\phi_\pi}\\
&=\frac{\kappa(\phi_\pi-1)+(1-\beta)\phi_y}{\sigma+\phi_y+\kappa\phi_\pi}>0
\end{tightalign}
\[
\Rightarrow \kappa(\phi_\pi-1)+(1-\beta)\phi_y>0\quad \blue{\text{satisfy if}}.
\]
\[
3^\circ\quad
1+\operatorname{tr}(A_T)+\det(A_T)
=\frac{(\sigma+\phi_y+\kappa\phi_\pi)+(\sigma+\kappa+\beta(\sigma+\phi_y))+\beta\sigma}{\sigma+\phi_y+\kappa\phi_\pi}>0
\quad \blue{\text{satisfy}}.
\]
\[
\red{\text{The system is stable iff: }}\quad \kappa(\phi_\pi-1)+(1-\beta)\phi_y>0.
\]
\noindent
\par\smallskip
\red{If $\phi_y=0\Rightarrow \phi_\pi>1$, consistent with the Taylor principle:
when inflation rises by $1\%$, the central bank needs to raise the nominal interest rate by $>1\%$.}
\par\smallskip
\[
r_t=i_t-\E_t\pi_{t+1}
\qquad \text{Fisher equation.}
\]

\[
\red{\E_t\pi_{t+1}\uparrow 1\%,\quad
i_t\uparrow>1\%
\Rightarrow
r_t\uparrow
\Rightarrow
\E_t\pi_{t+1}\downarrow.}
\]
\newpage
\section*{Monetary policy shock}
Assume $v_t$ follows AR(1):
\[
v_t=\rho_vv_{t-1}+\varepsilon_t^v,
\qquad \rho_v\in[0,1].
\]

\[
\varepsilon_t^v>0: \text{ contractionary MP shock},
\qquad \varepsilon_t^v<0: \text{ expansionary MP shock.}
\]

\par\smallskip
For simplicity, shut down technology shock:
\[
\hat a_t=0\Rightarrow \hat y_t^n=0,
\quad \hat r_t^n=0,\quad \forall t.
\]
Guess
\[
\tilde y_t=\psi_{yv}v_t,
\qquad \hat\pi_t=\psi_{\pi v}v_t.
\]
Plug into (***):
\begin{tightalign}
\psi_{yv}v_t&=-\frac{1}{\sigma}(\phi_\pi\psi_{\pi v}v_t+\phi_y\psi_{yv}v_t+v_t-\rho_v\psi_{\pi v}v_t)+\rho_v\psi_{yv}v_t,\\
-\sigma\psi_{yv}&=\phi_\pi\psi_{\pi v}+\phi_y\psi_{yv}+1-\rho_v\psi_{\pi v}-\sigma\rho_v\psi_{yv},\\
[\sigma(1-\rho_v)+\phi_y]\psi_{yv}&=(\rho_v-\phi_\pi)\psi_{\pi v}-1\quad (1).
\end{tightalign}
From NKPC:
\begin{tightalign}
\hat\pi_t&=\beta\E_t\{\hat\pi_{t+1}\}+\kappa\tilde y_t,\\
\psi_{\pi v}v_t&=\beta\rho_v\psi_{\pi v}v_t+\kappa\psi_{yv}v_t,\\
\psi_{yv}&=\frac{1-\beta\rho_v}{\kappa}\psi_{\pi v}\quad (2).
\end{tightalign}
Plug (2) into (1):
\[
\psi_{\pi v}=-\frac{\kappa}{(1-\beta\rho_v)[\sigma(1-\rho_v)+\phi_y]+\kappa(\phi_\pi-\rho_v)}.
\]
\[
\hat\pi_t
=
\underbrace{-\kappa\Lambda_\nu}_{<0}
v_t
\]
\[
\Lambda_v=\frac{1}{(1-\beta\rho_v)[\sigma(1-\rho_v)+\phi_y]+\kappa(\phi_\pi-\rho_v)}>0.
\]
where
\[
\kappa
=
\left(\sigma+\frac{\varphi+\alpha}{1-\alpha}\right)\lambda
=
\left(\sigma+\frac{\varphi+\alpha}{1-\alpha}\right)
\frac{(1-\theta)(1-\beta\theta)}{\theta}
\frac{1-\alpha}{1-\alpha+\alpha\varepsilon}
>0.
\]
\[
\tilde y_t
=
\underbrace{-(1-\beta\rho_v)\Lambda_v}_{<0}
v_t
\]
\begin{tightalign}
\hat r_t
&=
\hat i_t-\E_t\{\hat\pi_{t+1}\}
=
\phi_\pi\psi_{\pi v}v_t
+
\phi_y\psi_{yv}v_t
+
v_t
+
\kappa\Lambda_v\rho_vv_t
\\
&=
\underbrace{\sigma(1-\rho_v)(1-\beta\rho_v)\Lambda_v}_{>0}
v_t.
\end{tightalign}
\[
\hat i_t
=
\underbrace{\left[\sigma(1-\rho_v)(1-\beta\rho_v)-\kappa\rho_v\right]\Lambda_v}_{\text{coefficient on }v_t}
v_t.
\]
\[
\text{ambiguous; depends on }\rho_v.\quad \text{if }\rho_v\uparrow,\ \hat i_t\downarrow.
\]
\blue{If monetary policy is more persistent, $\hat i_t\downarrow$ real interest rate decreases.}
\[
\text{Therefore}
\left\{
\begin{aligned}
v_t &= \rho_v v_{t-1}+\varepsilon_t^v,
\qquad \rho_v\in[0,1],\\
\hat\pi_t &= -\kappa\Lambda_v v_t,\\
\tilde y_t &= -(1-\beta\rho_v)\Lambda_v v_t,\\
\hat r_t &= \sigma(1-\rho_v)(1-\beta\rho_v)\Lambda_v v_t,\\
\hat i_t &= \left[\sigma(1-\rho_v)(1-\beta\rho_v)-\kappa\rho_v\right]\Lambda_v v_t.
\end{aligned}
\right.
\]

\par\smallskip
\noindent
\blue{Calibration: $\beta=0.99$ (quarterly data), $\sigma=1$ (log utility), $\varphi=1$, $\alpha=1/3$, $\varepsilon=6$, $\theta=2/3$, $\bar y=4$.}

\noindent
\blue{With $\phi_\pi=1.5$, $\phi_y=0.5$, $\rho_v=0.5$, and shock $\varepsilon_t^v=25$bp: 
$\hat\pi_t<0$, $\tilde y_t<0$, $\hat y_t<0$, $\hat i_t>0$, $\hat m_t<0$.}
\section*{Technology shock}
Assume an AR(1) process of $\hat a_t$:
\[
\hat a_t=\rho_a\hat a_{t-1}+\varepsilon_t^a,
\qquad \rho_a\in[0,1].
\]
\[
\hat r_t^n=\sigma\E_t\{\hat y_{t+1}^n-\hat y_t^n\}=\sigma\psi_{ya}^n\E_t\{\Delta\hat a_{t+1}\}
=-\sigma\psi_{ya}^n(1-\rho_a)\hat a_t.
\]
Setting $v_t=0,\forall t$. Assume:
\[
\tilde y_t=\psi_{yr}\hat r_t^n,
\qquad \hat\pi_t=\psi_{\pi r}\hat r_t^n.
\]
\[
\begin{bmatrix}\tilde y_t\\\hat\pi_t\end{bmatrix}
=A_T\begin{bmatrix}\E_t\{\tilde y_{t+1}\}\\\E_t\{\hat\pi_{t+1}\}\end{bmatrix}
+B_T(\hat r_t^n-v_t).
\]
Opposite sign to $v_t$:
\[
\psi_{yr}=-\psi_{yv},\qquad \psi_{\pi r}=-\psi_{\pi v}.
\]
\begin{tightalign}
\tilde y_t&=(1-\beta\rho_a)\Lambda_a\hat r_t^n
=-\sigma\psi_{ya}^n(1-\rho_a)(1-\beta\rho_a)\Lambda_a\hat a_t,\\
\hat\pi_t&=\kappa\Lambda_a\hat r_t^n=-\sigma\psi_{ya}^n(1-\rho_a)\kappa\Lambda_a\hat a_t,
\end{tightalign}
where
\[
\Lambda_a=\frac{1}{(1-\beta\rho_a)[\sigma(1-\rho_a)+\phi_y]+\kappa(\phi_\pi-\rho_a)}>0.
\]

\[
\begin{aligned}
\hat y_t
&=
\hat y_t^n+\tilde y_t
=
\psi_{ya}^n\hat a_t
-
\sigma\psi_{ya}^n(1-\rho_a)(1-\beta\rho_a)\Lambda_a\hat a_t
\\
&=
\psi_{ya}^n
\left[
1-
\underbrace{\sigma(1-\rho_a)(1-\beta\rho_a)\Lambda_a}_{(*)}
\right]\hat a_t.
\end{aligned}
\]

\[
\begin{aligned}
1-(*) 
&=
1-
\sigma(1-\rho_a)(1-\beta\rho_a)
\frac{1}{
\sigma(1-\rho_a)(1-\beta\rho_a)
+
\phi_y(1-\beta\rho_a)
+
\kappa(\phi_\pi-\rho_a)
}
\\
&=
\frac{
\sigma(1-\rho_a)(1-\beta\rho_a)
+
\phi_y(1-\beta\rho_a)
+
\kappa(\phi_\pi-\rho_a)
-
\sigma(1-\rho_a)(1-\beta\rho_a)
}{
\sigma(1-\rho_a)(1-\beta\rho_a)
+
\phi_y(1-\beta\rho_a)
+
\kappa(\phi_\pi-\rho_a)
}
\\
&=
\frac{
\phi_y(1-\beta\rho_a)
+
\kappa(\phi_\pi-\rho_a)
}{
\sigma(1-\rho_a)(1-\beta\rho_a)
+
\phi_y(1-\beta\rho_a)
+
\kappa(\phi_\pi-\rho_a)
}.
\end{aligned}
\]

\[
\Rightarrow
\hat y_t
=
\frac{
\phi_y(1-\beta\rho_a)
+
\kappa(\phi_\pi-\rho_a)
}{
\sigma(1-\rho_a)(1-\beta\rho_a)
+
\phi_y(1-\beta\rho_a)
+
\kappa(\phi_\pi-\rho_a)
}
\psi_{ya}^n\hat a_t.
\]


\[
\begin{aligned}
\phi_y(1-\beta\rho_a)+\kappa(\phi_\pi-\rho_a)
&=
\kappa(\phi_\pi-1)
+
(1-\beta)\phi_y
+
\kappa
-
\rho_a\kappa
+
\beta\phi_y
-
\beta\rho_a\phi_y
\\
&=
\underbrace{\kappa(\phi_\pi-1)+(1-\beta)\phi_y}_{\text{stable condition}}
+
\underbrace{(1-\rho_a)\kappa}_{0\le \rho_a<1}
+
\underbrace{(1-\rho_a)\beta\phi_y}_{0\le \rho_a<1}
>0.
\end{aligned}
\]

\blue{$y_t$ moves positively with $a_t$.}
\[
(1-\alpha)\hat n_t=\hat y_t-
\hat a_t
=\left[(\psi_{ya}^n-1)-\sigma\psi_{ya}^n(1-\rho_a)(1-\beta\rho_a)\Lambda_a\right]\hat a_t.
\]
In calibration: $\sigma=1\Rightarrow \psi_{ya}^n=\frac{1+\varphi}{\varphi+\alpha+\sigma(1-\alpha)}=1\Rightarrow n_t$ declines when $a_t$ increases.


\par\smallskip
\section*{Equilibrium with exogenous money supply: $\Delta \hat m_t$}
\[
\frac{M_t}{P_t}=L(i_t,Y_t).
\]
\red{Close the model with money demand equation instead of Taylor rule.}
Define log real money balance:
\[
\widehat{\mathbb{M}}_t
=
\hat m_t-\hat p_t
\Rightarrow
\widehat{\mathbb{M}}_t
=
\hat y_t-\eta \hat i_t.
\]

\[
\Rightarrow
\tilde y_t-\eta\hat i_t
=
\widehat{\mathbb{M}}_t-\hat y_t^n
\Rightarrow
\hat i_t
=
\frac{1}{\eta}
\left(
\tilde y_t+\hat y_t^n-\widehat{\mathbb{M}}_t
\right).
\]

\[
\red{\text{Output}\uparrow
\Rightarrow
\text{money demand}\uparrow
\Rightarrow
i\uparrow.
\quad
\text{Real money supply}\uparrow
\Rightarrow
i\downarrow.}
\]

\[
DIS
\Rightarrow
(1+\sigma\eta)\tilde y_t
=
\sigma\eta\E_t\{\tilde y_{t+1}\}
+
\widehat{\mathbb{M}}_t
+
\eta\E_t\{\hat\pi_{t+1}\}
+
\eta\hat r_t^n
-
\hat y_t^n.
\]

From $\widehat{\mathbb{M}}_t$ definition:
\[
\widehat{\mathbb{M}}_t
-
\widehat{\mathbb{M}}_{t-1}
=
\Delta\hat m_t-\hat\pi_t
\Rightarrow
\widehat{\mathbb{M}}_{t-1}
=
\widehat{\mathbb{M}}_t+\hat\pi_t-\Delta\hat m_t.
\]

\[
A_{M,0}
\begin{bmatrix}
\tilde y_t\\
\hat\pi_t\\
\widehat{\mathbb{M}}_{t-1}
\end{bmatrix}
=
A_{M,1}
\begin{bmatrix}
\E_t\{\tilde y_{t+1}\}\\
\E_t\{\hat\pi_{t+1}\}\\
\widehat{\mathbb{M}}_t
\end{bmatrix}
+
B_M
\begin{bmatrix}
\hat r_t^n\\
\hat y_t^n\\
\Delta\hat m_t
\end{bmatrix}.
\]
$\widehat{\mathbb{M}}_{t-1}$ is predetermined and $\tilde y_t$, $\hat\pi_t$ are not predetermined. Stationary solution will exist iff
\[
A_M=A_{M,0}^{-1}A_{M,1}\text{ has two eigenvalues inside the unit circle.}
\]
Assume the monetary policy shock follows AR(1):
\[
\Delta\hat m_t=\rho_m\Delta\hat m_{t-1}+\varepsilon_t^m,
\qquad \rho_m\in[0,1].
\]
Let
\[
\hat r_t^n=\hat y_t^n=0,
\quad \varepsilon_t^m=0.25,
\quad t=0,
\quad \varepsilon_t^m=0,
\quad \forall t\ne0.
\]
An expansionary monetary policy shock:
\[
\Delta\hat m_t\uparrow\Rightarrow \tilde y_t\uparrow,\ \hat\pi_t\uparrow,\ \widehat{\mathbb{M}}_t\uparrow,
\qquad \hat i_t\uparrow:\text{ no liquidity effect }(\sigma=1),
\qquad \hat r_t\downarrow.
\]
\[
\hat a_t\uparrow\Rightarrow \hat r_t^n\downarrow,\ \hat y_t\uparrow,\ \hat\pi_t\downarrow,\ \hat n_t\downarrow,\ \hat r_t\uparrow,\ \hat i_t\approx0,
\qquad \hat y_t^n\uparrow\uparrow,\ \tilde y_t\downarrow.
\]
\blue{Liquidity effect: when real money balance $\uparrow$, $\widehat{\mathbb{M}}_t\uparrow\Rightarrow$ households hold too much money than they want.}
\[
\blue{\Rightarrow \hat i_t\downarrow\text{ to induce households to hold that money.}}
\]
\red{Why $\sigma=1$ liquidity effect is gone?}
\[
\hat i_t
=
\frac{1}{\eta}
\left(
\tilde y_t+\hat y_t^{\,n}-\widehat{\mathbb{M}}_t
\right)
\quad
\text{with } \hat y_t^{\,n}=0.
\]
\red{DIS:}
\[
\tilde y_t=-\frac{1}{\sigma}(\hat i_t-
\E_t\hat\pi_{t+1})+
\E_t\tilde y_{t+1}.
\]
Set $\hat y_t^n=0$:
\[
\hat i_t=\frac{1}{\eta}(\tilde y_t-
\widehat{\mathbb{M}}_t).
\]
\red{Liquidity effect: $\Delta\hat m_t>0$, $\hat i_t<0$, requires $\widehat{\mathbb{M}}_t>\tilde y_t$.}
\red{Real balance must rise by more than the output gap.}
\[
\red{
DIS \Rightarrow \hat i_t
=
\underbrace{\E_t\hat\pi_{t+1}}_{\red{>0}}
+
\sigma
\underbrace{\left[\E_t\hat y_{t+1}-\hat y_t\right]}_{\red{<0}}.
}
\]
\[
\hat i_t<0\Leftrightarrow \sigma(\tilde y_t-\E_t\tilde y_{t+1})>\E_t\hat\pi_{t+1}.
\]

\par\smallskip
\red{After expansionary money shock, $\tilde y_t\uparrow\Rightarrow \tilde y_t>\E_t\tilde y_{t+1}\Rightarrow \E_t\tilde y_{t+1}-\tilde y_t<0$.}

\par\smallskip
\[
\red{\text{When }\sigma=1,\text{ the negative term is not large enough to dominate the positive term.}}
\]

\par\smallskip
\blue{Liquidity effect disappears whenever:}
\[
\E_t\hat\pi_{t+1}>\tilde y_t-
\E_t\tilde y_{t+1}.
\]
Thus, $\sigma\gg1$, $\rho_m\ll1$ \blue{liquidity effect presents; money shock not persistent $\Rightarrow$ expected inflation $\downarrow$.}

\ThisNoteEnd
